\section{Introduction}
\label{sec:intro}

Trading is increasingly delegated to algorithms that learn from their own
profits. A growing literature shows that such algorithms can arrive at
\emph{supracompetitive} outcomes without communicating: reinforcement-learning
(RL) pricing agents settle on prices above the competitive level
\citep{calvano2020,klein2021}, algorithmic market makers quote wider spreads
\citep{colliard2026,cont2024}, and RL informed speculators trade less
aggressively than competition would imply, earning collusive profits and
degrading price informativeness \citep{dou2025}. For regulators the finding is
alarming, because tacit collusion by algorithms falls outside the reach of
rules written around agreements among humans \citep{harrington2018,ezrachi2016}.

A supracompetitive outcome, however, is not the same thing as collusion. In the
game-theoretic sense, collusion is an outcome sustained by the threat of
punishment: a trader who deviates to grab a larger share of the rent is
punished by its rivals, so deviating does not pay. The same low-intensity
outcome can instead come from a learning bias: an algorithm that systematically
undervalues aggressive actions will trade conservatively even when no rival
could ever punish it. \citet{dou2025} name both channels in the trading context
(``price-trigger'' collusion and ``over-pruning'' bias), and a parallel debate
in industrial organization asks whether Q-learning collusion is genuine or an
artifact of exploration and learning design
\citep{calvano2023,abada2023,banchio2023,asker2022,denboer2026}. The
distinction is not academic. If algorithms punish, transparency matters,
because monitoring is what makes punishment possible, and the natural remedies
target information: reporting delays, reduced order-flow disclosure. If the
outcome is a learning bias, transparency is irrelevant and the remedies target
algorithm design.

This paper asks how to tell the two apart, and what the answer is for RL
informed traders in the canonical \citet{kyle1985} market and in the
\citet{dou2025} extension with information-insensitive investors.

\paragraph{Contributions.}
\begin{enumerate}
\item \textbf{Diagnostics that can fail, and do.} We combine a paired
  \emph{rival-deviation} test (one agent is forced to play its one-period
  best response; we measure the rivals' reaction under common random numbers)
  and the \citet{dou2025} \emph{noise-shock} test with placebos that cannot
  collude: a \emph{myopic} learner ($\gamma=0$), \emph{no memory}, an
  \emph{uninformative memory} of the same size as an informative one, and
  \emph{passive traders} that move the collusive outcome above the Nash one,
  so that collusion and bias predict opposite directions. Applied to the
  logit Bertrand game of \citet{calvano2020}, the diagnostics find punishment
  (the rival cuts its price and deviating loses), and they flag the same
  game without memory, where prices are even higher ($\Delta=\BertNone$), as a
  supracompetitive outcome with no punishment behind it.
\item \textbf{When can traders punish at all?} In the standard Kyle market a
  one-period deviation from the collusive intensity moves order flow by only
  $(I-1)/(2I)$ noise standard deviations per unit of $v/\sigma_v$, whatever
  the volume of noise trading; with a large mass of information-insensitive
  investors the ratio is about 625 in the \citet{dou2025} calibration
  (Proposition~\ref{prop:snr}). Building on \citet{green1984}, we show that
  in the standard market Nash-reversion trigger strategies cannot sustain
  \emph{any} collusion, so a near-cartel outcome would need much harsher
  punishments, and we bound what any punishment scheme can sustain
  (Propositions~\ref{prop:trigger} and~\ref{prop:anypun}). A myopic learner
  whose estimates are unbiased plays Nash (Proposition~\ref{prop:myopic}), so
  the myopic placebo measures learning bias directly.
\item \textbf{The bias, isolated.} A single informed trader with no rival and
  a fixed price rule under-trades by up to 57\% under standard Q-learning, by
  an amount proportional to $\sqrt\alpha$ ($R^2=\MechSqrtRsq$), while
  counterfactual updating learns the optimum.
\item \textbf{Low trading without collusion.} In the standard Kyle market,
  forward-looking Q-learners close on average \ResidDeltaInt{} of the gap
  between competition and the cartel, yet rivals never punish deviations,
  even with perfect monitoring, and deviating pays. The myopic placebo trades
  even less ($\Delta=\MyopicResidDeltaInt$, beyond the cartel); an uninformative
  memory reproduces the outcome ($\RandomMemDeltaInt$); and when passive
  traders put the cartel \emph{above} Nash, learners still trade half their
  Nash intensity and earn \PassiveProfitPct\% less than in equilibrium.
  Blurring the order flow changes nothing ($\Delta=\OpaqueDeltaInt$), while
  counterfactual updating, a change to the learning rule alone, restores
  competition ($\Delta=\CfResidDeltaInt$).
\item \textbf{Where deviations are glaring.} In the \citet{dou2025} regime,
  forward-looking learners still never react to shocks or deviations, while a
  \emph{myopic} shared-table learner reacts to both: it shows the
  threshold-shaped response to price shocks that has been read as a price
  trigger, although it has no reason to punish and deviating still pays.
  These signatures do not identify collusion without a placebo.
\item \textbf{An audit protocol and an open testbed.} We assemble the
  diagnostics into a five-step protocol that returns \emph{collusion} for the
  \citet{calvano2020} baseline and \emph{learning bias} for every market of
  informed traders we study, and apply it to competing Q-learning dealers
  under adverse selection. We release a batched, CPU-only simulator with
  exact benchmarks, tabular Q-learning, DQN and PPO agents, a compiled
  training engine, all diagnostics, and a test suite that validates the
  economics and the diagnostics against scripted agents whose behaviour is
  known.
\end{enumerate}

\paragraph{Why this matters for policy.} The two channels call for different
interventions. Our results suggest that in markets like the standard Kyle
setting, the supracompetitive outcomes produced by off-the-shelf Q-learning are
properties of the learning procedure, and would respond to algorithm-design
standards (step sizes, exploration) rather than to limits on transparency.
Where deviations are highly detectable, as with a large mass of
information-insensitive investors, the diagnostics must be applied with a
placebo, because a reaction to a price shock is not by itself evidence of
punishment.
